\chapter{Station C: Darf ich das posten?}\label{ch:posten}

\begin{myremark}
	Das Meme von der schlafenden Lena finden viele lustig -- Lena nicht.
	Durfte das überhaupt jemand posten? An dieser Station lernen Sie die wichtigsten Regeln für Bilder, Videos und Musik kennen.
\end{myremark}

\section{Zwei Rechte, die Sie kennen müssen}

\begin{mysummary}[Urheberrecht]
	Wer etwas erschafft -- ein Foto, eine Zeichnung, einen Text, ein Lied --, ist \tib{Urheberin} oder \tib{Urheber} und entscheidet, ob und wie das Werk veröffentlicht wird.
	Geregelt ist das im \gls{URG}.
	\begin{itemize}
		\item In der Schweiz ist seit 2020 \textbf{jedes Foto} geschützt, auch ein einfacher Schnappschuss.
		\item Im \textbf{privaten Kreis} (Familie, enge Freunde) dürfen Sie Werke nutzen und teilen.
		      Öffentlich posten gehört \emph{nicht} mehr zum privaten Kreis.
		\item Ein Bild aus der Internetsuche ist darum \emph{nicht} automatisch frei verwendbar.
	\end{itemize}
\end{mysummary}

\begin{mysummary}[Recht am eigenen Bild]
	Jede Person darf grundsätzlich selbst bestimmen, ob Bilder von ihr veröffentlicht werden.
	Das folgt aus dem Persönlichkeitsschutz im \gls{ZGB}.
	\begin{itemize}
		\item Wer auf einem Bild \textbf{erkennbar} ist und im Mittelpunkt steht, muss in der Regel \textbf{einverstanden} sein.
		\item Bei einer grossen Menschenmenge (Konzert, Sportanlass) ist das meist nicht nötig, solange niemand gezielt herausgehoben wird.
		\item Blossstellende oder verletzende Bilder sind besonders heikel -- sie können Cybermobbing sein.
	\end{itemize}
\end{mysummary}

\begin{myexercise}[Fallkarten: Darf ich das posten?]\label{ex:fallkarten}
	Entscheiden Sie zu zweit: \faCheck{} erlaubt, \faTimes{} nicht erlaubt oder \textbf{?} kommt darauf an. Begründen Sie kurz.

	\begin{center}
		\renewcommand{\arraystretch}{1.5}
		\small
		\begin{tabular}{p{0.6cm}p{7.4cm}cp{3.8cm}}
			\textbf{Nr.} & \textbf{Situation}                                                                                        & \textbf{\faCheck/\faTimes/?} & \textbf{Begründung} \\
			\midrule
			1            & Ein Selfie, auf dem nur Sie zu sehen sind, als Profilbild.                                                &                              &                     \\
			2            & Ein Gruppenfoto der Klasse auf Ihrem öffentlichen Profil, ohne die anderen zu fragen.                       &                              &                     \\
			3            & Das Foto der schlafenden Lena als Meme im Klassenchat.                                                    &                              &                     \\
			4            & Ein Bild aus der Internet-Bildersuche als Titelbild für Ihren Blog.                                       &                              &                     \\
			5            & Ein Bild von Wikimedia Commons mit der Lizenz \gls{CC}~BY auf einem Schulplakat, mit Angabe der Urheberin. &                              &                     \\
			6            & Ihr Lieblingssong als Hintergrundmusik in einem eigenen Kurzvideo.                                        &                              &                     \\
			7            & Ein Foto vom Konzert mit Tausenden Zuschauenden.                                                          &                              &                     \\
			8            & Ein Screenshot aus einem privaten Chat mit einer Kollegin in Ihrer Story.                                 &                              &                     \\
		\end{tabular}
	\end{center}

	\begin{myanswer}
		\begin{enumerate}
			\item \faCheck{} Sie sind Urheber und abgebildete Person zugleich.
			\item \faTimes{} Alle erkennbaren Personen müssen einverstanden sein.
			\item \faTimes{} Verletzt Lenas Recht am eigenen Bild und ist blossstellend (Cybermobbing). Auch ein Klassenchat ist nicht mehr nur der enge private Kreis.
			\item \faTimes{} Bilder aus der Suche sind in der Regel urheberrechtlich geschützt. Lizenz prüfen oder freie Bilder suchen.
			\item \faCheck{} Die Lizenz erlaubt die Nutzung, wenn die Urheberin korrekt genannt wird.
			\item \textbf{?} Kommt auf die Plattform an: Manche haben Lizenzen für ihre eigene Musikbibliothek. Einen Song selbst hochzuladen und öffentlich zu verwenden, ist ohne Erlaubnis nicht zulässig.
			\item \faCheck{} Menschenmenge, niemand steht im Mittelpunkt. Das Foto haben Sie selbst gemacht.
			\item \faTimes{} Private Nachrichten sind vertraulich; die Kollegin muss einverstanden sein.
		\end{enumerate}
	\end{myanswer}
\end{myexercise}

\section{Freie Bilder finden und richtig angeben}

Mit einer \tib{Creative-Commons-Lizenz} (\gls{CC}) erlauben Urheberinnen und Urheber allen, ihr Werk unter bestimmten Bedingungen zu nutzen:

\begin{center}
	\renewcommand{\arraystretch}{1.3}
	\begin{tabular}{lp{11cm}}
		\textbf{Kürzel} & \textbf{Bedeutung}                                                                 \\
		\midrule
		BY              & Namensnennung: Die Urheberin oder der Urheber muss genannt werden.                 \\
		SA              & Weitergabe unter gleichen Bedingungen: Veränderte Werke brauchen dieselbe Lizenz.  \\
		NC              & Nicht kommerziell: keine Nutzung, mit der Geld verdient wird.                      \\
		ND              & Keine Bearbeitung: Das Werk darf nicht verändert werden.                           \\
		CC0             & Kein Vorbehalt: frei nutzbar, Namensnennung freiwillig.                            \\
	\end{tabular}
\end{center}

Für eine korrekte Angabe hilft die Regel \tib{TULLU}: \textbf{T}itel, \textbf{U}rheber/in, \textbf{L}izenz, \textbf{L}ink zur Lizenz, \textbf{U}rsprungsort.

\begin{myexample}
	\enquote{Rotkehlchen im Schnee} von Max Muster, Lizenz CC~BY~4.0 (\url{https://creativecommons.org/licenses/by/4.0/}), Quelle: Wikimedia Commons.
\end{myexample}

\begin{myexercise}[Ein freies Bild für Lena]
	Lena möchte ein neues Profilbild ohne ihr Gesicht, zum Beispiel ein Tier oder eine Landschaft.
	Suchen Sie auf \url{https://openverse.org} oder \url{https://commons.wikimedia.org} ein passendes Bild mit einer \gls{CC}-Lizenz
	und schreiben Sie die vollständige Angabe nach TULLU auf.
	\schreibzeile
	\schreibzeile
\end{myexercise}

\begin{mychallenge}[Unter welcher Lizenz steht dieses Skript?]
	Finden Sie auf dem Titelblatt oder in der Kopfzeile heraus, unter welcher Lizenz dieses Skript steht
	(Tipp: Folgen Sie dem Link zu FreeFlower). Erklären Sie, was Sie damit tun dürfen und was nicht.
	\schreibzeile
	\schreibzeile
	\begin{myanswer}
		Das Skript steht unter \gls{CC}~BY-NC-SA~4.0: Man darf es kopieren, weitergeben und verändern,
		muss die Urheber nennen, darf damit kein Geld verdienen und muss veränderte Fassungen unter derselben Lizenz weitergeben.
	\end{myanswer}
\end{mychallenge}

\begin{mychallenge}[Und KI-Bilder?]
	Mit \gls{KI} lassen sich Bilder in Sekunden erzeugen. Diskutieren Sie zu zweit: Wem gehört ein solches Bild?
	Darf man ein KI-Bild einer echten Person posten, das diese Person nie so gesehen hat?
	Halten Sie Ihre wichtigsten Argumente fest. (Hinweis: Die Rechtslage ist noch nicht in allen Punkten geklärt. Klar ist aber, dass das Recht am eigenen Bild auch hier gilt.)
	\schreibzeile
	\schreibzeile
\end{mychallenge}

\begin{fallakte}[-- Station C]
	Durfte das Meme von Lena gepostet werden? Was kann Lena jetzt tun?
	\schreibzeile
	\schreibzeile
	\schreibzeile
	\begin{myanswer}
		Nein, Lena war erkennbar und nicht einverstanden; das Bild stellt sie bloss.
		Lena kann die Person bitten, das Bild zu löschen, es der Plattform bzw. im Chat melden, Screenshots als Beweis sichern
		und sich an eine Vertrauensperson (Eltern, Klassenlehrperson) wenden. Beratung bietet z.\,B. Pro Juventute (Telefon 147).
	\end{myanswer}
\end{fallakte}
